\documentclass[12pt,a4paper]{article}
\usepackage[margin=2.3cm]{geometry}
\usepackage[brazil]{babel}
\usepackage[utf8]{inputenc}
\usepackage[T1]{fontenc}
\usepackage{booktabs}
\usepackage{graphicx}
\usepackage{float}
\usepackage{amsmath}
\usepackage{hyperref}
\usepackage{caption}
\usepackage{array}
\usepackage{longtable}

\title{Evidências empíricas sobre Transformers para previsão e ordenação de ativos da B3}
\author{Relatório técnico preliminar}
\date{\today}

\begin{document}
\maketitle

\section{Objetivo e desenho empírico}

Os experimentos avaliam se previsões geradas por arquiteturas baseadas em Transformers contêm informação útil para a ordenação cross-section de ativos da B3. A análise é feita com métricas aderentes à decisão de carteira: IC de Spearman, precisão direcional positiva, precisão direcional negativa e AUC direcional.

A separação por base de dados é necessária. Preços, retornos simples e log-retornos representam problemas econométricos diferentes. A previsão de preços explora persistência e dinâmica do nível; a previsão de retornos simples e de log-retornos é mais próxima do problema de ranking financeiro usado na literatura de seleção de ações. Por isso, as Tabelas \ref{tab:precos}--\ref{tab:logret} não misturam bases.

Os benchmarks considerados são \textit{EqualWeight}, \textit{Momentum} e \textit{RandomG}. Nos artefatos exportados, \textit{RandomG} aparece como grupo \textit{Benchmark} com modelo \texttt{k\_1}, \texttt{k\_5}, etc. O relatório usa a nomenclatura econômica, não a nomenclatura interna do arquivo.

\section{Métricas}

O IC de Spearman é a correlação ordinal entre o score previsto e o retorno realizado:
\[
IC_t^{Sp}=\rho\left(\operatorname{rank}(\hat{s}_{i,t}),\operatorname{rank}(r_{i,t})\right).
\]
A precisão positiva mede a frequência com que os ativos selecionados para compra apresentam retorno positivo. A precisão negativa mede o análogo para o lado vendido. A AUC de alta e a AUC de queda medem a capacidade do score de separar, respectivamente, retornos positivos e negativos. Valores de AUC acima de $0,5$ indicam discriminação superior ao aleatório.

Em todas as tabelas, ``Med.'' é a mediana entre configurações disponíveis e ``Max.'' é o maior valor observado entre configurações comparáveis. Para benchmarks com apenas uma linha agregada nos artefatos locais, a coluna de máximo coincide com a estatística agregada disponível.

\section{Resultados em $k=1$ por base de dados}

\begin{table}[H]
\centering
\caption{Base de preços: estatísticas medianas e máximas em $k=1$.}
\label{tab:precos}
\resizebox{\textwidth}{!}{%
\begin{tabular}{lrrrrrrrrrrr}
\toprule
Modelo & n & IC Med. & IC Max. & Prec.+ Med. & Prec.+ Max. & Prec.- Med. & Prec.- Max. & AUC alta Med. & AUC alta Max. & AUC queda Med. & AUC queda Max. \\
\midrule
MASTER | MASTER & 3 & 0.102 & 0.143 & 51.0\% & 65.8\% & 47.4\% & 47.6\% & 0.548 & 0.593 & 0.569 & 0.591 \\
TFB | DUET & 3 & 0.090 & 0.091 & 53.3\% & 53.5\% & -- & -- & 0.533 & 0.534 & 0.533 & 0.534 \\
TFB | FEDformer & 3 & 0.095 & 0.102 & 53.5\% & 53.8\% & -- & -- & 0.535 & 0.541 & 0.535 & 0.541 \\
TFB | Nonstationary\_Transformer & 3 & 0.097 & 0.102 & 53.9\% & 54.0\% & -- & -- & 0.531 & 0.534 & 0.531 & 0.534 \\
TFB | TimesNet & 3 & 0.093 & 0.098 & 53.2\% & 54.1\% & -- & -- & 0.532 & 0.534 & 0.532 & 0.533 \\
\bottomrule
\end{tabular}}
\end{table}

Na base de preços, o MASTER apresenta o maior IC máximo em $k=1$, com $0,143$ na agregação por configuração. Entre os modelos TFB, FEDformer e Nonstationary Transformer têm ICs máximos ligeiramente acima de $0,10$, e TimesNet apresenta a maior precisão positiva máxima entre os TFB, com $54,1\%$.

\begin{table}[H]
\centering
\caption{Base de retornos simples: estatísticas medianas e máximas em $k=1$.}
\label{tab:ret}
\resizebox{\textwidth}{!}{%
\begin{tabular}{lrrrrrrrrrrr}
\toprule
Modelo & n & IC Med. & IC Max. & Prec.+ Med. & Prec.+ Max. & Prec.- Med. & Prec.- Max. & AUC alta Med. & AUC alta Max. & AUC queda Med. & AUC queda Max. \\
\midrule
MASTER | MASTER & 3 & 0.017 & 0.046 & 46.8\% & 47.9\% & 46.2\% & 49.7\% & 0.511 & 0.517 & 0.512 & 0.513 \\
TFB | DUET & 3 & 0.051 & 0.065 & 51.9\% & 52.2\% & -- & -- & 0.522 & 0.527 & 0.524 & 0.528 \\
TFB | FEDformer & 3 & 0.021 & 0.024 & 50.1\% & 50.2\% & -- & -- & 0.507 & 0.507 & 0.508 & 0.508 \\
TFB | Nonstationary\_Transformer & 3 & 0.016 & 0.037 & 49.3\% & 50.6\% & -- & -- & 0.508 & 0.517 & 0.507 & 0.518 \\
TFB | TimesNet & 3 & 0.030 & 0.054 & 51.0\% & 51.3\% & -- & -- & 0.513 & 0.523 & 0.515 & 0.524 \\
\bottomrule
\end{tabular}}
\end{table}

Na base de retornos simples, o sinal é menor, mas permanece positivo para os melhores modelos TFB. O DUET apresenta IC máximo de $0,065$ e precisão positiva máxima de $52,2\%$. Essa evidência é mais conservadora que a obtida em preços e deve ser tratada como a comparação mais próxima do problema de ranking de retornos.

\begin{table}[H]
\centering
\caption{Base de log-retornos: estatísticas medianas e máximas em $k=1$.}
\label{tab:logret}
\resizebox{\textwidth}{!}{%
\begin{tabular}{lrrrrrrrrrrr}
\toprule
Modelo & n & IC Med. & IC Max. & Prec.+ Med. & Prec.+ Max. & Prec.- Med. & Prec.- Max. & AUC alta Med. & AUC alta Max. & AUC queda Med. & AUC queda Max. \\
\midrule
MASTER | MASTER & 3 & -0.004 & 0.004 & 49.2\% & 51.9\% & 46.1\% & 46.2\% & 0.496 & 0.497 & 0.504 & 0.506 \\
TFB | DUET & 3 & 0.050 & 0.065 & 52.0\% & 52.2\% & -- & -- & 0.522 & 0.527 & 0.523 & 0.528 \\
TFB | FEDformer & 3 & 0.022 & 0.026 & 50.2\% & 50.4\% & -- & -- & 0.508 & 0.508 & 0.508 & 0.509 \\
TFB | Nonstationary\_Transformer & 3 & 0.016 & 0.037 & 49.4\% & 50.6\% & -- & -- & 0.508 & 0.517 & 0.507 & 0.518 \\
TFB | TimesNet & 3 & 0.031 & 0.055 & 51.2\% & 51.4\% & -- & -- & 0.514 & 0.523 & 0.515 & 0.524 \\
\bottomrule
\end{tabular}}
\end{table}

Os log-retornos reproduzem a hierarquia observada nos retornos simples: DUET tem o melhor desempenho entre os TFB, seguido por TimesNet. A interpretação é que há informação ordinal útil, mas a magnitude do sinal em retornos é inferior à observada em preços.

\begin{table}[H]
\centering
\caption{Benchmarks em $k=1$: RandomG, Momentum e EqualWeight.}
\label{tab:bench}
\resizebox{\textwidth}{!}{%
\begin{tabular}{lrrrrrrrrrr}
\toprule
Benchmark & IC Med. & IC Max. & Prec.+ Med. & Prec.+ Max. & Prec.- Med. & Prec.- Max. & AUC alta Med. & AUC alta Max. & AUC queda Med. & AUC queda Max. \\
\midrule
RandomG & -0.015 & -0.015 & 48.8\% & 48.8\% & -- & -- & 0.494 & 0.494 & 0.494 & 0.494 \\
Momentum & -0.015 & -0.015 & 48.6\% & 48.6\% & -- & -- & 0.494 & 0.494 & 0.494 & 0.494 \\
EqualWeight & -- & -- & 49.2\% & 49.2\% & -- & -- & -- & -- & -- & -- \\
\bottomrule
\end{tabular}}
\end{table}

Os benchmarks têm função de controle. Momentum e RandomG ficam abaixo dos melhores Transformers em IC, precisão positiva e AUC. EqualWeight não gera ranking cross-section propriamente dito; por isso, IC e AUC não são definidos nos artefatos exportados, embora a precisão direcional da carteira agregada possa ser calculada.

\section{Precisão dos modelos TFB por horizonte}

A Tabela \ref{tab:tfbprec} explicita a precisão positiva dos modelos TFB por \texttt{pred\_len} e janela de trading. Esta tabela responde diretamente ao ponto da taxa de acerto direcional: a precisão média sai de valores próximos de $51\%$ em $k=1$ para valores próximos ou superiores a $60\%$ quando a janela de realização aumenta.

\begin{table}[H]
\centering
\caption{TFB: precisão positiva por \texttt{pred\_len} e janela de trading.}
\label{tab:tfbprec}
\resizebox{\textwidth}{!}{%
\begin{tabular}{rrrrrrrr}
\toprule
pred\_len & k & n & Prec. média & Prec. med. & Prec. max. & Sharpe médio & Retorno médio \\
\midrule
1 & 1 & 36 & 51.0\% & 50.8\% & 54.2\% & 2.868 & 74.7\% \\
5 & 1 & 36 & 50.8\% & 50.2\% & 54.4\% & 2.740 & 72.2\% \\
5 & 5 & 36 & 55.7\% & 56.0\% & 62.4\% & 2.988 & 75.1\% \\
10 & 1 & 36 & 51.1\% & 51.0\% & 54.5\% & 2.950 & 76.2\% \\
10 & 5 & 36 & 56.2\% & 56.7\% & 63.2\% & 3.194 & 79.0\% \\
10 & 10 & 36 & 59.3\% & 59.9\% & 67.6\% & 2.798 & 76.5\% \\
15 & 1 & 36 & 51.0\% & 51.1\% & 54.7\% & 2.987 & 77.2\% \\
15 & 5 & 36 & 55.8\% & 56.4\% & 63.3\% & 3.123 & 77.1\% \\
15 & 10 & 36 & 59.6\% & 60.4\% & 68.3\% & 2.877 & 79.2\% \\
15 & 15 & 36 & 61.7\% & 64.3\% & 72.5\% & 2.837 & 79.8\% \\
24 & 1 & 36 & 50.9\% & 50.9\% & 54.1\% & 2.868 & 72.8\% \\
24 & 5 & 36 & 55.7\% & 56.8\% & 62.4\% & 2.861 & 71.2\% \\
24 & 10 & 36 & 59.1\% & 59.7\% & 67.3\% & 2.669 & 74.8\% \\
24 & 15 & 36 & 62.0\% & 64.1\% & 72.4\% & 2.933 & 82.3\% \\
24 & 20 & 36 & 62.0\% & 63.9\% & 79.8\% & 2.715 & 77.4\% \\
24 & 24 & 36 & 62.3\% & 64.0\% & 72.9\% & 1.954 & 67.8\% \\
\bottomrule
\end{tabular}}
\end{table}

\begin{figure}[H]
    \centering
    \includegraphics[width=0.82\textwidth]{figuras/fig_tfb_precision_line_predlen.png}
    \caption{TFB: precisão positiva média por \texttt{pred\_len} e janela de trading.}
    \label{fig:tfbprecmean}
\end{figure}

\begin{figure}[H]
    \centering
    \includegraphics[width=0.82\textwidth]{figuras/fig_tfb_precision_max_line_predlen.png}
    \caption{TFB: precisão positiva máxima por \texttt{pred\_len} e janela de trading.}
    \label{fig:tfbprecmax}
\end{figure}

A leitura técnica das Figuras \ref{fig:tfbprecmean} e \ref{fig:tfbprecmax} é que a acurácia direcional não cresce por um único ponto isolado da grade. Há deslocamento sistemático da precisão para cima quando a janela de trading aumenta. Essa regularidade sugere que parte do sinal capturado pelos Transformers se materializa em janelas de realização mais longas, e não necessariamente no pregão imediatamente seguinte.

\section{Gráficos facetados por modelo}

Os gráficos facetados por modelo são os mais adequados para discutir heterogeneidade entre arquiteturas. Com os artefatos disponíveis localmente, foi possível reconstituir os gráficos facetados por modelo em função da janela de trading $k$ para Spearman e precisão. Esses gráficos entram como diagnóstico complementar nas Figuras \ref{fig:spearkfacet} e \ref{fig:preckfacet}.

\begin{figure}[H]
    \centering
    \includegraphics[width=\textwidth]{figuras/fig_global_spearman_por_k_facet_modelo.png}
    \caption{IC Spearman mediano por janela de trading, facetado por modelo.}
    \label{fig:spearkfacet}
\end{figure}

\begin{figure}[H]
    \centering
    \includegraphics[width=\textwidth]{figuras/fig_global_precision_por_k_facet_modelo.png}
    \caption{Precisão positiva mediana por janela de trading, facetada por modelo.}
    \label{fig:preckfacet}
\end{figure}

A figura exata pedida para o relatório final é a facetada por modelo com \texttt{pred\_len} no eixo horizontal e, alternadamente, IC Spearman e precisão positiva no eixo vertical. Essa figura depende do arquivo longo \texttt{comparativo\_metricas\_long\_com\_bench\_neg.csv}, que não está nos artefatos locais nem versionado no repositório. O relatório, portanto, deve ser atualizado com essa figura assim que esse CSV for disponibilizado ou quando o notebook TFB for reexecutado no ambiente onde o arquivo existe. A seção não inclui código para preservar o formato científico do texto.

\section{MASTER versus TFB}

A comparação entre MASTER e TFB exige cautela. Nos modelos TFB, \texttt{pred\_len} representa uma trajetória prevista de múltiplos passos. Na implementação do MASTER usada nos experimentos, o modelo prediz um ponto $H$ passos à frente. Assim, o MASTER não produz um vetor intermediário entre $t+1$ e $t+H$, o que impede a construção de certos intervalos de predição baseados na variabilidade ao longo do vetor previsto.

\section{Conclusão}

A evidência empírica é favorável à factibilidade dos Transformers, desde que a análise seja separada por base. Em preços, a habilidade de ranking em $k=1$ é forte. Em retornos simples e log-retornos, o sinal é menor, porém positivo para os melhores modelos TFB. A evidência mais robusta está na dinâmica por horizonte: as precisões direcionais crescem de forma relevante quando a janela de trading aumenta, alcançando valores máximos superiores a $70\%$ em algumas regiões da grade experimental.

\begin{thebibliography}{9}
\bibitem{Kwiatkowski2025} Kwiatkowski, J.; Chudziak, J. A. (2025). \textit{On Evaluating Loss Functions for Stock Ranking: An Empirical Analysis with Transformer Model}. arXiv:2510.14156.
\bibitem{Li2023MASTER} Li, T.; Liu, Z.; Shen, Y.; Wang, X.; Chen, H.; Huang, S. (2023). \textit{MASTER: Market-Guided Stock Transformer for Stock Price Forecasting}. arXiv:2312.15235.
\bibitem{Qiu2024TFB} Qiu, X. et al. (2024). \textit{TFB: Towards Comprehensive and Fair Benchmarking of Time Series Forecasting Methods}. Proceedings of the VLDB Endowment.
\bibitem{Fawcett2006} Fawcett, T. (2006). An introduction to ROC analysis. \textit{Pattern Recognition Letters}, 27(8), 861--874.
\end{thebibliography}

\end{document}
